% !TeX program = xelatex
% ============================================================
%  CHƯƠNG 5: THIẾT KẾ ĐIỀU PHỐI VÀ THỰC THI (PHASE 3 & 4)
%  PAD-ONAP: Proactive AI-Driven DDoS Defense on ONAP
% ============================================================
\documentclass[12pt,a4paper]{article}

\usepackage[a4paper,margin=2.5cm]{geometry}
\usepackage{fontspec}
\usepackage{polyglossia}
\setmainlanguage{vietnamese}
\setotherlanguage{english}
\setmainfont{Times New Roman}

\usepackage{amsmath,amssymb}
\usepackage{graphicx}
\usepackage{booktabs}
\usepackage{longtable}
\usepackage{array}
\usepackage{float}
\usepackage{enumitem}
\usepackage{listings}
\usepackage{xcolor}
\usepackage[colorlinks=true,linkcolor=blue,citecolor=blue,urlcolor=blue]{hyperref}

\setlength{\parindent}{1.2em}
\setlength{\parskip}{0.5em}

\lstset{
  basicstyle=\ttfamily\footnotesize,
  backgroundcolor=\color{gray!8},
  frame=single,
  breaklines=true,
  columns=fullflexible,
  keepspaces=true,
}

\title{\textbf{Chương 5: Thiết Kế Điều Phối và Thực Thi NFV}\\[0.4em]
\large PAD-ONAP: Proactive AI-Driven DDoS Defense trên kiến trúc ONAP\\[0.3em]
\normalsize Phase 3 \& 4 --- Orchestration và Service Function Chaining}
\author{}
\date{}

\begin{document}
\maketitle
\tableofcontents
\newpage

% ============================================================
\section{Tổng Quan Giải Pháp M3/M4}
% ============================================================

Sau khi mô hình Trí tuệ Nhân tạo ở cụm M2 (Phase 2) phát hiện hoặc dự báo được nguy cơ tấn công DDoS, hệ thống cần một "Bộ não Điều phối" để chuyển dịch những con số phi tuyến tính thành các lệnh thực thi cấu hình mạng vật lý và ảo hóa. Khối chức năng M3 (Orchestration) và M4 (Enforcement) được thiết kế dựa trên tiêu chuẩn kiến trúc ETSI MANO và ONAP (Open Network Automation Platform).

Kiến trúc tổng thể của việc điều phối được thực hiện như sau:
\begin{enumerate}[leftmargin=2em]
\item \textbf{Tier Mapper (M3a):} Thu nhận \texttt{AIOutputPayload} thông qua bus tin nhắn DMaaP/Kafka để quy đổi Confidence và Dự báo $P_{30s}$ sang Thang đo rủi ro 5 tầng (\textit{5-Tier Response}).
\item \textbf{Policy Engine (M3b):} Cổng gác (Guard) xử lý độ trễ (Hysteresis) và duy trì tần số (Frequency Guard) để kích hoạt tự động hóa vòng lặp kín (Closed-Loop Automation) dựa trên CLAMP.
\item \textbf{SLA Allocator (M3c):} Điều tiết băng thông theo thời gian thực để duy trì Cam kết chất lượng dịch vụ (SLA) cho các Mạng con (Network Slices).
\item \textbf{ONAP SO \& SFC (M4):} Khởi tạo và thiết lập các phần tử mạng ảo (VNF) thông qua Docker/Kubernetes và điều hướng lưu lượng (Traffic Steering) qua OVS (Open vSwitch).
\end{enumerate}

% ============================================================
\section{Thiết Kế M3 --- Policy Orchestration}
% ============================================================

\subsection{Cơ chế Tier Mapper và Tự động VNF}

Dựa trên tín hiệu Output của XGBoost (tốc độ cao) và Transformer+LSTM (dự báo dài hạn), hệ thống chia cách thức phản ứng ra thành 5 mức độ (Tier):

\begin{table}[H]
\centering
\caption{Ánh xạ Quyết định phòng thủ (Tier Mapping)}
\label{tab:tier_mapping}
\begin{tabular}{lp{6cm}p{6cm}}
\toprule
\textbf{Tier} & \textbf{Luật Điều kiện} & \textbf{Hành động ONAP (Action)} \\
\midrule
\textbf{T0} & Traffic Bình thường. & Không hành động. \\
\textbf{T1} & Conf $> 0.50$ & \textbf{Low:} Alert và Rate-limit dòng TCP/UDP. \\
\textbf{T2} & $P_{30s} > 0.70$ hoặc Conf $> 0.80$ & \textbf{Pre-position:} Khởi động VNF Firewall phòng hờ. \\
\textbf{T3} & $P_{30s} > 0.90$ và Conf $> 0.90$ & \textbf{Scale-out:} Kéo traffic vào Scrubber VNF. \\
\textbf{T4} & Lặp liên tục ở T3. & \textbf{Critical:} Cô lập hoàn toàn (Blackhole). \\
\bottomrule
\end{tabular}
\end{table}

\subsection{Luật Policy Engine: Tránh chập chờn hệ thống}

Trong thực tế, khi lưu lượng dao động ngay quanh ngưỡng báo động, hệ thống mạng sẽ phải đối mặt với \textit{Ping-Pong Effect} (bật và tắt VNF liên tục, làm quá tải tài nguyên Controller và gián đoạn Service). Cụm M3 giải quyết bài toán này bằng hai module thiết yếu:

\begin{enumerate}[leftmargin=2em]
\item \textbf{Hysteresis Constraint (Ngưỡng lưu trễ):}
Yêu cầu $N=2$ cửa sổ (\textit{windows}) liên tiếp có dấu hiệu tấn công thì mới nâng điểm \textit{Escalate} (đưa từ $T0 \rightarrow T1$). Ngược lại, yêu cầu $M=5$ cửa sổ liên tiếp sạch để \textit{De-escalate} gỡ bỏ VNF phòng thủ.
\item \textbf{Frequency Guard (Lá chắn tần suất):}
Khoảng thời gian nghỉ biến thiên \texttt{MIN\_INTERVAL\_S = 10.0s}. Nếu có lệnh chuyển đổi VNF, hệ thống khoá không nhận lệnh mới trong vòng 10 giây bất chấp các giá trị Confidence tăng vọt. Mọi lệnh trong khoảng thời gian này sẽ chuyển sang trạng thái \texttt{HOLD}.
\end{enumerate}

\begin{lstlisting}[language=json,caption={JSON Log Payload Cảnh báo Proactive với Policy Guard HOLD}]
{
  "event_id": "0e19fd9a-0f4f...",
  "attack_type": "Amplification",
  "tier": 1,
  "action": "NEW_ATTACK",
  "acted": true,
  "proactive": false
}
// Window ke tiep bi chan boi Frequency Guard
{
  "event_id": "5377de53-3ed0...",
  "attack_type": "Amplification",
  "tier": 1,
  "action": "HOLD",
  "acted": false
}
\end{lstlisting}

\subsection{SLA Allocator --- Bài toán tái cấu trúc QoS}

Băng thông (Total bandwidth) cấp phát cho Datacenter là hữu hạn. Khi đẩy VNF Scrubber vào trong OVS, \textit{Overhead} (độ hao hụt) ăn mòn vào tổng băng thông thực.
Hệ thống sử dụng bài toán \textbf{Weighted Fair Queuing Allocation}:
Giả lập 3 Network Slices mạng 5G:
\begin{itemize}
    \item \texttt{slice-eMBB} (Weight: 1.0, Guaranteed: 200 Mbps)
    \item \texttt{slice-URLLC} (Weight: 2.0, Guaranteed: 100 Mbps)
    \item \texttt{slice-mMTC} (Weight: 0.5, Guaranteed: 50 Mbps)
\end{itemize}
Thuật toán đảm bảo phân bổ lại băng thông một cách công bằng ngay cả khi VNF Overhead chiếm dụng lên đến hàng chục Mbps khi ở Tier 3 / Tier 4.

% ============================================================
\section{Thiết Kế M4 --- Thực Thi NFV và SFC}
% ============================================================

\subsection{Truy vấn ONAP Service Orchestrator (SO)}

Các phần tử bảo mật VNF (Firewall, Scrubber, Analyzer) được định trước trong SDNC catalog qua YAML Helm Chart / Docker Images. Chức năng \texttt{onap\_so\_client.py} đóng vai trò gọi webhook gRPC/REST xuống VIM (Virtual Infrastructure Manager) để sinh tài nguyên cô lập namespace.

Dòng đời (Lifecycle): `Instantiate $\rightarrow$ Health Check Loop (Wait Active) $\rightarrow$ SFC Steering $\rightarrow$ Terminate`.

Lưu ý: Môi trường hiện tại phục vụ mục đích đồ án đánh giá vòng lặp kín AI, nên API SO được triển khai dưới dạng Stub (chưa điều phối Kubernetes vật lý thật), trong đó giả lập \texttt{TimeoutError} độ trễ tương tự khi boot instance trong thực tế (1-3s).

\subsection{Service Function Chaining (SFC)}
Đóng vai trò định tuyến luồng tin dị thường di chuyển ngang (service chain) sang VNF. 
Logic SFC Manager: Nếu $Tier \ge 2$, chèn OpenFlow rule vào Switch ảo chặn nguồn Mac/Port rác và đẩy qua Mac giả lập của VNF Scrubber. Việc này thực thi gần tốc độ dòng ($O(1)$) trên Ternary Content-Addressable Memory (TCAM).

% ============================================================
\section{Đánh Giá Thực Nghiệm Hệ Thống Toàn Trình (P1$\rightarrow$P4)}
% ============================================================

Để chứng minh luận điểm, hệ thống PAD-ONAP hoàn chỉnh đã được khởi tạo chạy nền và chịu tải tiêm thử nghiệm (Anomaly Injector) dạng UDP Flood.

\subsection{Độ trễ toàn trình (E2E Latency Tracker)}

Độ trễ luôn là một trong những điểm yếu của NFV khi tích hợp Trí tuệ Nhân tạo Cấp cao. PAD-ONAP tính toán và gửi dữ liệu xuất \textit{Prometheus Exporter} với port \texttt{9292}. Hệ thống theo dõi bốn Latch quan trọng để bóc tách:

\begin{table}[H]
\centering
\caption{Kết quả đo E2E Latency trên Sample Logs}
\label{tab:latency_results}
\begin{tabular}{lrrr}
\toprule
\textbf{Công đoạn điều phối} & \textbf{Thời gian P50} & \textbf{P95} & \textbf{P99} \\
\midrule
AI $\rightarrow$ Policy (\texttt{d2p\_ms})     & $<1.00$ ms & $1.52$ ms & $2.20$ ms \\
Policy $\rightarrow$ SO (\texttt{p2s\_ms})       & $0.41$ ms  & $0.85$ ms & $1.15$ ms \\
SO $\rightarrow$ VNF Active (\texttt{s2v\_ms})   & $102$ ms   & $145$ ms  & $210$ ms (*) \\
VNF $\rightarrow$ SFC Steer (\texttt{v2s\_ms})   & $2.14$ ms  & $3.50$ ms & $4.10$ ms \\
\midrule
\textbf{Tổng Toàn Trình (Total E2E)} & \textbf{$\approx 105$ ms} & \textbf{$\approx 150$ ms} & \textbf{$\approx 217$ ms}\\
\bottomrule
\multicolumn{4}{l}{\footnotesize (*) Thời gian kéo container giả lập và health-check. Khác biệt với độ trễ inference.}
\end{tabular}
\end{table}

\textbf{Đánh giá:} Mô hình AI (Phase 2) có độ trễ $<10$ ms. Công đoạn lâu nhất nằm ở việc hệ thống gọi lệnh SO kéo VNF lên (thường trên lý thuyết tốn vài giây). Quyết định "Pre-position VNF từ sớm" (Proactive T+30s / Tier 2) do Track B cung cấp khiến thời gian khởi động này bị ẩn đi (Hidden Overhead), nhờ vậy traffic chưa từng bị drop ngay khi cuộc tấn công đạt điểm rơi cao nhất. Hệ thống vận hành \textit{Near Real-time} hoàn hảo.

\subsection{Tổng Kết Đóng Góp Công Trình}

Tựu chung lại, Bốn Phase triển khai đã thể hiện trọn vẹn ưu điểm thiết kế của PAD-ONAP:
\begin{itemize}
    \item Lớp \textbf{Telemetry} gNMI gọn nhẹ, bóc tách 17 cấu trúc mạng.
    \item Lớp \textbf{AI Explainable \& Proactive} khắc phục nhược điểm "Cái Hộp Đen" của Machine Learning, giúp dự báo 30s với độ chuẩn xác $>96\%$.
    \item Lớp \textbf{Orchestration \& Guard} triệt tiêu "Ping-Pong Effect", điều dưỡng băng thông SLA ổn định ngay trong bão tấn công.
\end{itemize}

% ============================================================
\end{document}
